\documentclass[11pt,a4paper]{article}
\usepackage[margin=24mm]{geometry}
\usepackage[T1]{fontenc}
\usepackage{lmodern,microtype,amsmath,amssymb,booktabs,array,enumitem,xurl,longtable}
\usepackage{xr-hyper}
\usepackage{hyperref}
\externaldocument[R-]{../output/pdf/blowup_density_revised}[blowup_density_revised.pdf]
\hypersetup{colorlinks=true,linkcolor=blue,urlcolor=blue,citecolor=blue,
 pdftitle={Guide to this Project}}
\newcommand{\code}[1]{\nolinkurl{#1}}
\newcommand{\paperref}[1]{\ref{R-#1}}
\newcommand{\source}[4]{\code{#1.lean:#2}, #4 \code{#3}}
\newcommand{\mapped}[2]{#1~\paperref{#2}}
\newcommand{\coverage}[1]{\textbf{#1.} }
\setlength{\emergencystretch}{2em}
\setlist{nosep}
\title{Guide to this Project}
\author{}
\date{21 September 2026}
\begin{document}
\maketitle
\noindent This guide accompanies the article. It maps
its numbered statements to proof locations, compares the cited and implemented
proof methods, and indexes reusable classical results. The article's
Section~\paperref{sec:preliminaries} fixes the analytic framework;
Sections~\paperref{sec:torus} and~\paperref{sec:whole} treat the torus and
whole space. Source locations refer to the accompanying repository snapshot.

\section{Build requirements}\label{guide:build}
Use Lean \code{4.34.0-rc2} and the checked-in manifests, which pin Mathlib to
\code{85e3a25e006c35636f0e53b0e9296caca2685bc0}. Build and test in Linux:
\begin{verbatim}
lake -d verification exe cache get
make check
make test
make test-mutations
\end{verbatim}
The suite checks 30 current publication interfaces. The separate declaration
audit checks the transitive axioms of the results cited in this guide.
For macOS 27, use a Linux VM because of the pinned runtime's allocator issue.
The repository README gives the build details.

\section{Finding the theorem corresponding to the article}\label{guide:correspondence}
Each entry gives the proof file, declaration line and theorem name.
The prefixes \code{F/} and \code{B/} mean
\code{formalization/NSFormalization/} and \code{verification/Bindings/},
respectively. Names are local to the indicated file's namespace; line numbers
refer to the source tree identified above. A \code{def} entry is explicitly a constructed
witness, not a theorem declaration. Some final proofs assemble several
results in \code{Bindings}; these are distinct from the conformance wrappers
in \code{verification/Tests/}.

\textbf{Closed} means that the complete article statement has a Lean
kernel-checked proof, including every auxiliary result and its instantiation.
No unproved theorem input, \code{sorry}/\code{admit}, or extra axiom is allowed.
Only the article's stated hypotheses and the standard logical axioms
\code{propext}, \code{Classical.choice}, and \code{Quot.sound} remain.
\textbf{Partial} means that only part of the article statement is formally proved.
A downstream result can be Closed when it uses only fully formalized cases of a Partial result.

All 26 numbered statements and the numbered remark appear below.
The table follows the article's numbering; only its two main theorems retain
parenthetical titles. Qualifications
record differences in scope. Historical identifiers retain \emph{insertion}
and \emph{packet} where the article uses \emph{gluing} and \emph{building block}.

\begingroup\small
\setlength{\tabcolsep}{3pt}
\setlength{\LTpre}{6pt}
\renewcommand{\arraystretch}{1.15}
\begin{longtable}{@{}>{\raggedright\arraybackslash}p{0.20\linewidth}>{\raggedright\arraybackslash}p{0.76\linewidth}@{}}
\toprule Article statement & Proof location and scope \\\midrule
\endfirsthead
\toprule Article statement & Proof location and scope (continued) \\\midrule
\endhead
\bottomrule\endfoot
\mapped{Theorem}{thm:packet} &
\coverage{Closed}\source{F/Source/PacketBreakdown}{76}{source_breakdown}{theorem}.
Construction and same-force global nonexistence for every positive viscosity,
with blowup time fixed at one. \\
\mapped{Proposition}{prop:local} &
\coverage{Partial}\source{B/LocalTheoryV2}{93}{localTheoryV2}{def};
\source{F/Section4/A02/Uniqueness}{120}{velocity_unique}{theorem};
\source{F/Section4/A02/Maximal}{157}{exists_maximal_of_localSolution}{theorem};
\source{F/Section4/A04/ShiftedExtension}{247}{extendsBeyond_of_memForceR'}{theorem};
\source{B/TorusLocalTheory}{483}{torusLocalTheory}{def}.
Proved for the density force classes, with $H^7/H^3$ restart.
The general $H^1$-uniform restart statements are not proved
(Section~\ref{guide:methods}). \\
\mapped{Lemma}{lem:packetenergy} & \coverage{Closed}\source{F/Section3/T14/PacketEnergy}{160}{energy_le_work_of_packet}{theorem}; \source{F/Section3/T14/PacketEnergy}{106}{work_eq_square_of_packet}{theorem}; \source{F/Section4/I01/Quiet}{53}{exists_packet_quiet}{theorem}. \\
\mapped{Theorem}{thm:main} & \coverage{Closed}\source{F/Section3/T21/MainAssembly}{84}{mainStatement_holds}{theorem}. \\
\mapped{Lemma}{lem:localization} & \coverage{Closed}\source{F/Section3/T13/Assembly}{330}{localizationAPI}{theorem}. Includes the interior-order estimate and the two endpoint identities. \\
\mapped{Proposition}{prop:scaling} & \coverage{Closed}\source{F/Section3/T15/Assembly}{130}{scalingStatement_holds}{theorem}. \\
\mapped{Lemma}{lem:potential} & \coverage{Closed}\source{F/Section3/T16/Assembly}{483}{localPotential}{theorem}. \\
\mapped{Lemma}{lem:correction} & \coverage{Closed}\source{F/Section3/T17/ArticleScope}{44}{correctionStatementArticle_holds}{theorem};
\source{F/Section3/T17/ArticleScope}{70}{article_force_identification}{theorem};
\source{F/Section3/T17/ArticleScope}{119}{article_correction_derivative_bound}{theorem};
\source{F/Section3/T17/ArticleScope}{129}{article_force_derivative_bound}{theorem};
\source{F/Section3/T17/ArticleScope}{103}{article_force_spatial_volume}{theorem};
\source{F/Section3/T17/ArticleScope}{111}{article_force_time_length}{theorem};
\source{F/Section3/T17/ArticleScope}{140}{article_correction_energy_bound}{theorem};
\source{F/Section3/T17/ArticleScope}{155}{article_force_mixed_bound}{theorem};
\source{F/Section3/T17/ArticleScope}{173}{article_force_sobolev_bound}{theorem}.
Classical reference, building-block placement and the Lemma 3.4 ball; zero extension agrees on the classical slab and preserves every displayed bound. \\
\mapped{Theorem}{thm:insertion} & \coverage{Closed}\source{F/Section3/T19/FromData}{20}{periodicInsertion_from_data}{theorem};
\source{B/PeriodicInsertionV2}{10}{periodicInsertionStatementV2_holds}{theorem}.
From raw data for every prescribed coordinate ball with closure in the cube interior;
all four clauses, including shrinking support in that chart, hold simultaneously. \\
\mapped{Proposition}{prop:density} & \coverage{Closed}\source{F/Section3/T19/Assembly}{43}{periodicDensityStatement_holds}{theorem}. \\
\mapped{Proposition}{prop:critical} & \coverage{Closed}\source{F/Section3/T20/Assembly}{53}{criticalRegularityStatement_holds}{theorem}. \\
\mapped{Corollary}{cor:nondensity} & \coverage{Closed}\source{F/Section3/T21/MainAssembly}{80}{nonDensityStatement_unconditional}{theorem}. \\
\mapped{Corollary}{cor:mixed} & \coverage{Closed}\source{F/Section3/T19/Assembly}{47}{mixedRegionStatement_holds}{theorem}. \\
\mapped{Corollary}{cor:closure} & \coverage{Closed}\source{F/Section3/T19/Assembly}{51}{strongClosureStatement_holds}{theorem}. \\
\mapped{Proposition}{prop:projection} & \coverage{Closed}\source{F/Section3/T19/Assembly}{55}{projectionStatement_holds}{theorem}. \\
\mapped{Remark}{rem:peaks} & \coverage{Closed}\source{F/Section3/T18/ForceAmplitude}{24}{packetForce_ne_zero}{theorem};
\source{F/Section3/T18/ForceAmplitude}{240}{forceAmplitude_lower}{theorem};
\source{F/Section3/T18/ForceAmplitude}{288}{forceAmplitude_diverges}{theorem};
\source{F/Section3/T18/ForceAmplitude}{332}{forceAmplitude_real_diverges}{theorem};
\source{B/ForceAmplitude}{40}{forceAmplitude_from_data}{theorem}.
The positive finite packet amplitude gives the displayed lower bound and divergence for every inserted family. \\
\mapped{Corollary}{cor:boundary} & \coverage{Closed}\source{B/BoundaryInsertionV2}{27}{boundaryInsertion_from_data}{theorem};
\source{F/Section3/T22/Assembly}{25}{boundedDomainNorm}{theorem}.
\\
\mapped{Proposition}{prop:affine} & \coverage{Closed}\source{B/AffineVariation}{109}{affineVariationPacket}{theorem}. Instantiated at the selected upstream building block. \\
\mapped{Proposition}{prop:multiple} & \coverage{Partial}\source{F/Section3/T24/MultipleAssembly}{51}{multipleRegionsStatement_holds}{theorem}.
Periodic case; the bounded-domain variant is not included. \\
\mapped{Proposition}{prop:conservative} & \coverage{Partial}\source{F/Section3/T24/ConservativeAssembly}{45}{conservativeForcing}{theorem}.
Periodic case; the bounded-domain variant is not included. \\
\mapped{Theorem}{thm:Rmain} & \coverage{Closed}\source{B/MainThresholds}{85}{mainThresholds}{theorem}. Density: \source{B/DensityFromInsertion}{30}{breakdownDenseR_of_subcritical}{theorem}; converse: \source{B/MainThresholds}{34}{mainThresholds_nonDensity}{theorem}. \\
\mapped{Theorem}{thm:Rinsert} & \coverage{Closed}\source{B/InsertionFromData}{150}{wholeSpaceInsertion_holds}{theorem}.
Any prescribed nonempty open set, with one family and the given reference solution. \\
\mapped{Proposition}{prop:Rcritical1} & \coverage{Closed}\source{F/Section4/R43/Universal}{193}{universal_of_memForceR}{theorem}. \\
\mapped{Proposition}{prop:Rcritical2} & \coverage{Closed}\source{F/Section4/R44/Prop44}{24}{main}{theorem}.
Finite-horizon $L^2_tH^{-1/2}_x$ estimate in the inhomogeneous norm,
with zero initial velocity. \\
\mapped{Corollary}{cor:Rclasses} & \coverage{Closed}\source{B/ForceClasses}{46}{forceClasses}{theorem}. \\
\mapped{Proposition}{prop:Renergy} & \coverage{Closed}\source{B/CompletedDensity}{14}{completedDensity}{theorem}. \\
\mapped{Theorem}{thm:Rgrid} & \coverage{Closed}\source{B/GridObservations}{47}{gridObservations_choose}{theorem}.
Observation times are strictly below the terminal time. \\
\end{longtable}
\endgroup

\noindent\textbf{Upstream building block.}
Under \code{vendor/NavierStokesAndEuler/NavierStokes/},
\code{R3ActualCandidate.lean:18} proves
\code{NavierStokes.R3CompactCandidate.selected_compact_candidate};
\code{R3FiniteEnergyComparison.lean:37} proves
\code{NavierStokes.ComparatorBridge.compact_candidate_excludes_global_solution}.
In \code{ComparatorR3Theorem.lean:36},
\code{NavierStokes.ComparatorBridge.navier_stokes_breakdown_R3} combines the
candidate, uniqueness comparison and viscosity rescaling.
\code{Source/PacketBreakdown.lean} applies positive-viscosity uniqueness to the
selected construction and proves the complete statement of
Theorem~\paperref{thm:packet} for the same force.

\smallskip
\noindent The unnumbered analytic inputs are indexed with literature sources
and proof locations in Section~\ref{guide:classical}.

\section{Differences in proof methods}\label{guide:methods}
This section compares the proof routes used in the cited literature with the
routes implemented in the code. The coverage of each article statement is listed in
Section~\ref{guide:correspondence}.

\paragraph{Local theory and continuation (Proposition~\paperref{prop:local}).}
Tao's Theorems~5.1 and~5.4~\cite{tao} construct local solutions by a contraction
in the $H^1$ mild-solution space and then propagate higher regularity on the
same interval. The formalization instead uses its available $H^7$ local
construction on $\mathbb R^3$ and $H^3$ construction on the torus when restarting
a smooth solution. To obtain continuation from
$\int_0^S\|u(t)\|_{H^2}^2\,dt<\infty$, it first applies the high-order energy
estimate and Gr\"onwall:
\[
 \sup_{t<S}\|u(t)\|_{H^m}
 \le\left(\|a\|_{H^m}+\int_0^S\|f(t)\|_{H^m}\,dt\right)
 \exp\!\left(C_{m,\nu}\int_0^S\|u(t)\|_{H^2}^2\,dt\right)<\infty.
\]
Taking $m=7$ or $3$ gives a bounded set of restart data. For the fixed smooth
force on a compact time window, the code obtains one positive lifespan valid
for all these restart data and times, then pastes by uniqueness across $S$.
The general $H^1$-uniform restart statement is not proved; the continuation
argument above uses only the proved high-order case. Tao's higher-regularity argument
(Theorem~5.4(iv), p.~53) explains why smooth data remain in all Sobolev orders;
the uniform high-order restart bound is proved separately in the code.
See \source{F/Section4/A04/RestartFixedForce}{41}{restartFixedForce_of_memForceR}{theorem},
\source{F/Section4/A04/ShiftedExtension}{247}{extendsBeyond_of_memForceR'}{theorem}
and \source{F/Section3/T11/ExistenceInputH3}{456}{extendsBeyondH3}{theorem}.

\paragraph{Fractional localization (Lemma~\paperref{lem:localization}).}
The article invokes Taylor's coordinate-cover and partition-of-unity
construction of Sobolev spaces, with interpolation
\cite[Chapter~4, Sections~2--3]{taylor}. A fixed chart cutoff then gives the
required bounded local-to-global map. The formal proof establishes the
estimate directly for the torus: it identifies homogeneous fractional norms
with double-integral seminorms on $\mathbb R^3$ and on the torus, compares the
periodized kernel with the Euclidean kernel, and bounds the remaining lattice
tail by a multiple of $\|z\|_2^2$. The fixed distance from the support ball to
the chart boundary makes the constant uniform as the support shrinks.
The endpoint orders $0$ and $1$ are treated separately.
The kernel comparison is
\source{F/Section3/T13/KernelComparison}{394}{iTorus_periodize_le}{theorem};
the resulting estimate is
\source{F/Section3/T13/Assembly}{330}{localizationAPI}{theorem}. Includes the interior-order estimate and the two endpoint identities.

\paragraph{Critical regularity and non-density.}
The article invokes the forced critical fixed-point theory
\cite[Theorem~2.3.1 and Lemma~2.3.2, pp.~47--49]{danchin}, with the
inhomogeneous heat estimate (Theorem~2.2.5, p.~44) for the finite-horizon
$L^2_tH^{-1/2}_x$ case. The code instead proves the critical energy
inequalities, closes a smallness bootstrap and uses the high-order
continuation route above. See the proof locations for
Propositions~\paperref{prop:critical}, \paperref{prop:Rcritical1}
and~\paperref{prop:Rcritical2} in Section~\ref{guide:correspondence}. The open-ball argument yielding
non-density is unchanged; only the analytic proof route differs.

\section{Classical results available for future reference}\label{guide:classical}
The following classical results have formal proofs in this project and may
provide starting points for later formalizations. The entries give actual
proof locations using the prefixes from Section~\ref{guide:correspondence}. Each entry gives a literature locator and identifies any specialization or
derivation needed to reach the formal version. Unless stated otherwise, the spatial domain is
$\mathbb R^3$ and vector fields take values in $\mathbb R^3$.

\begin{enumerate}[leftmargin=*,itemsep=5pt]
\item \textbf{Sobolev embedding into bounded functions.}
For smooth fields with square-integrable derivatives through order two,
the code proves the pointwise and essential-supremum $H^2$ bounds:
\source{F/Section4/A03/BoundedRepresentative}{165}{enorm_le_jetENorm}{theorem}
and \source{F/Section4/A03/BoundedRepresentative}{182}{eLpNormTop_le_jetENorm}{theorem}.
These use the derivative-based norm \code{jetENorm}.
\emph{Source:} Tao~\cite[Appendix~A, (A.13), draft p.~336]{taodispersive},
with $d=3$, $s=2$. Apply the scalar estimate componentwise and use Plancherel
to pass to the equivalent derivative norm; the literal constant is not claimed
to coincide with the code's constant.

\item \textbf{Sobolev algebra and tame product estimates.}
At integer orders $m\ge2$, scalar products satisfy the $H^m$ algebra estimate,
and scalar--vector products satisfy the tame estimate with an $H^2$ low norm:
\source{F/Section4/A03/ScalarTameProduct}{308}{algebraProductScalar}{theorem}
and \source{F/Section4/A03/VectorTameProduct}{273}{tameProductVector}{theorem}.
\emph{Source:} Tao~\cite[Appendix~A, Lemma~A.8, (A.17)--(A.18),
draft p.~338]{taodispersive}. Formula~(A.18) gives the algebra estimate;
(A.17) combined with (A.13) gives the $H^2$ low norm in dimension three.
The vector version follows componentwise.

\item \textbf{Integral representations of fractional Sobolev norms.}
For $0<s<1$, the double-integral seminorm equals an explicit constant times
the squared homogeneous Fourier norm:
\source{F/Section3/T13/WholeSpaceIdentity}{509}{wholeSpace_identity}{theorem}
for smooth compactly supported fields, and
\source{F/Section3/T13/TorusIdentity}{1174}{torus_identity}{theorem}
for smooth periodic fields, with the mean removed in the homogeneous norm.
\emph{Sources:} Di~Nezza--Palatucci--Valdinoci~\cite[Proposition~3.4,
arXiv v3, pp.~16--17]{dpv} for $\mathbb R^3$; Roncal--Stinga~\cite[arXiv v3, Theorem~1.5,
(1.6)--(1.8), p.~3; Fourier definition (1.1), p.~1]{roncal} for the torus.
The periodic energy identity is a consequence of their operator formula:
set $\sigma=2s$, pair with the field and symmetrize the integral. Parseval
then gives the squared Fourier norm. Rescale their $2\pi$-periodic torus
to period one and adjust Fourier constants; the constant mode contributes zero.

\item \textbf{Smooth approximation in Sobolev and Bochner spaces.}
For every real $s$, compactly supported smooth vector fields are dense in the
implemented $H^s$ datum space:
\source{F/Section4/B01/Spatial}{134}{spatialApprox}{theorem}.
For $1\le q<\infty$, Sobolev-valued $L^q(0,\infty)$ paths can be approximated
by finite sums $\sum_j\phi_j(t)A_j$, with $\phi_j\in C_c^\infty(0,\infty)$:
\source{F/Section4/B01/Temporal}{170}{temporalApprox}{theorem}.
The temporal theorem alone does not make the spatial data $A_j$ smooth;
the spatial approximation supplies that separate step.
\emph{Sources:}~Taylor~\cite[Chapter~4, Section~1, Exercise~1,
author-chapter p.~5]{taylor} gives Schwartz density for every real $s$.
Multiplying a Schwartz approximant by $\chi(x/R)$ gives convergence in the
Schwartz topology and hence in $H^s$, yielding compact support; this last
step is a deduction, not the exercise's literal statement.
Hunter~\cite[Appendix~6.A, Proposition~6.29, p.~200]{hunter} gives the finite
separated-sum approximation on $(0,T)$. Truncating the integrable $q$-th
power tail first extends it to $(0,\infty)$ for $q<\infty$.

\item \textbf{A local vector potential for a divergence-free field.}
On a ball, a smooth time-dependent divergence-free field has a smooth vector
potential given by the radial integral:
\source{F/Section3/T16/BallPotential}{324}{exists_potential_on_ball}{theorem}.
\emph{Source:} Petersen~\cite[Section~7.1, Lemma~7.1.2, pp.~127--128,
and Lemma~7.1.5, p.~129]{petersen}. Apply the homotopy formula to the closed
$2$-form $\iota_v(dx_1\wedge dx_2\wedge dx_3)$ and the radial contraction
$x\mapsto x_0+\rho(x-x_0)$. Translating its primitive into vector notation
gives the code's radial integral. Smooth dependence on time follows by
differentiation under the integral over $0\le\rho\le1$. This specializes
the Poincar\'e lemma to a ball, rather than arbitrary domains.

\item \textbf{Classical Navier--Stokes local theory and continuation.}
Periodic local existence is proved in
\source{F/Section3/T11/ExistenceInputH3}{376}{exists_periodicLocalSolution_unconditional}{theorem}.
Whole-space velocity uniqueness and pressure uniqueness up to a time-dependent
constant are in
\source{F/Section4/A02/Uniqueness}{120}{velocity_unique}{theorem}
and \source{F/Section4/A02/Uniqueness}{244}{pressure_gauge}{theorem}.
The continuation theorems are listed in Section~\ref{guide:methods}. These results use the
project's smooth initial-data and force classes.
\emph{Sources:} Tao~\cite[Theorems~5.1(ii)--(iv), p.~48, and
5.4(ii)--(iv), pp.~52--53]{tao} for local existence, uniqueness and smoothness;
Lemma~4.1(i), (40), pp.~44--45, for the whole-space pressure gauge;
Corollary~5.2, p.~50, for periodic maximal development and its $H^1$ blowup
alternative. These are the classical local-theory inputs. The precise
squared-$H^2$ continuation test and fixed-force $H^7/H^3$ restart used here
are derived by the high-order estimate and Gr\"onwall in Section~\ref{guide:methods}; they are
not quoted as literal conclusions of Tao's theorems.
\end{enumerate}

\begin{thebibliography}{99}\small\raggedright
\setlength{\itemsep}{2pt}
\setlength{\parsep}{0pt}
\setlength{\parskip}{0pt}

\bibitem{tao}
T. Tao. Localisation and compactness properties of the Navier--Stokes global regularity problem. \emph{Analysis \& PDE} \textbf{6}(1) (2013), 25--107. \href{https://doi.org/10.2140/apde.2013.6.25}{doi:10.2140/apde.2013.6.25}.

\bibitem{taylor}
M. E. Taylor. \emph{Partial Differential Equations I: Basic Theory}. 2nd ed., Applied Mathematical Sciences, vol.~115. Springer, 2011. \href{https://doi.org/10.1007/978-1-4419-7055-8}{doi:10.1007/978-1-4419-7055-8}.

\bibitem{taodispersive}
T. Tao. \emph{Nonlinear Dispersive Equations: Local and Global Analysis}. CBMS Regional Conference Series in Mathematics, vol.~106. American Mathematical Society, 2006. \href{https://doi.org/10.1090/cbms/106}{doi:10.1090/cbms/106}.

\bibitem{dpv}
E. Di Nezza, G. Palatucci and E. Valdinoci. Hitchhiker's guide to the fractional Sobolev spaces. \emph{Bulletin des Sciences Math\'ematiques} \textbf{136}(5) (2012), 521--573. \href{https://doi.org/10.1016/j.bulsci.2011.12.004}{doi:10.1016/j.bulsci.2011.12.004}.

\bibitem{roncal}
L. Roncal and P. R. Stinga. Fractional Laplacian on the torus. \emph{Communications in Contemporary Mathematics} \textbf{18}(3) (2016), 1550033. \href{https://doi.org/10.1142/S0219199715500339}{doi:10.1142/S0219199715500339}.

\bibitem{hunter}
J. K. Hunter. \emph{Notes on Partial Differential Equations}. University of California, Davis, 2014. \url{https://www.math.ucdavis.edu/~hunter/pdes/pde_notes.pdf}.

\bibitem{petersen}
P. Petersen. \emph{Manifolds, Transversality, and de Rham Cohomology}. University of California, Los Angeles, lecture notes. \url{https://www.math.ucla.edu/~petersen/manifolds.pdf}.

\bibitem{danchin}
R. Danchin. \emph{Fourier Analysis Methods for PDE's}. Lecture notes, 2005. \url{https://perso.math.u-pem.fr/danchin.raphael/cours/courschine.pdf}.
\end{thebibliography}
\end{document}
